\subsection{抛物型方程的极值原理}\label{sec:parabolic-maximum-principles}

本节将最大值原理推广到如下形式的抛物型方程
\begin{equation}\label{eq:parabolic-maximum-principle-operator}
Lu=-\dfrac{\partial u}{\partial t}+a_{ij}(\bm{x},t)\dfrac{\partial^2u}{\partial x_i\partial x_j}+b_i(\bm{x},t)\dfrac{\partial u}{\partial x_i}+c(\bm{x},t)u=f(\bm{x},t).
\end{equation}


考虑定义在 $\Omega\times(0,T)$ 上的解, 其中 $\Omega$ 是 $\mathbb{R}^n$ 中的一个 (开) 区域 (也可以考虑 $(\bm{x},t)$ 空间中更一般的区域). 令 $D=\Omega\times(0,T]$, $Q=\Omega\times(0,T)$, 且 $\Sigma=(\partial\Omega\times[0,T])\cup(\Omega\times\{0\})$. 假设系数在 $\overline D$ 上连续, 且矩阵 $a_{ij}$ 在 $\overline D$ 上严格正定. 还始终假设 $u\in C^2(D)\cap C(\overline D)$. 抛物型方程的最大值原理指出, 在 $c$ 和 $f$ 满足适当符号条件时, $u$ 的最大值必在 $\Sigma$ 上.

\subsubsection{弱最大值原理}\label{subsec:parabolic-weak-maximum}

Theorem~\ref{thm:weak-maximum-principle} 对抛物型方程的对应结果如下.

\begin{Theorem}
\label{thm:parabolic-weak-maximum}
假设 $\Omega$ 有界且 $Lu\geq0$ ($Lu\leq0$). 再设在 $D$ 内 $c(\bm{x},t)=0$. 那么 $u$ 的最大值 (最小值) 在 $\Sigma$ 上达到.
\end{Theorem}

\begin{Proof}
由 Theorem~\ref{thm:weak-maximum-principle} 和 Remark~\ref{rem:weak-maximum-semidefinite}, 已知最大值在 $\partial D$ 上达到. 此外, 若 $Lu>0$, 则最大值不能在 $\Omega\times\{T\}$ 上达到. 事实上, 设 $\bm{x}\in\Omega\times\{T\}$ 满足 $u(\bm{x},T)=\max_{\bm{y}\in\Omega\times\{T\}}u(\bm{y},T)$. 那么在点 $\bm{x}$ 处有
\begin{equation}\label{eq:parabolic-top-spatial-maximum-inequality}
a_{ij}\dfrac{\partial^2u}{\partial x_i\partial x_j}+b_i\dfrac{\partial u}{\partial x_i}\leq0,
\end{equation}
由于 $Lu>0$, 得出 $u_t<0$.

对一般情形, 令 $u_\epsilon=u+\epsilon\exp(-t)$. 得到
\begin{equation}\label{eq:parabolic-weak-positive-perturbation}
Lu_\epsilon=Lu+\epsilon\exp(-t)>0,
\end{equation}
因而
\begin{equation}\label{eq:parabolic-weak-boundary-maximum}
\max_{\overline D}u_\epsilon=\max_\Sigma u_\epsilon.
\end{equation}
令 $\epsilon\to0$ 即得定理.
\end{Proof}

显然, Corollaries~\ref{cor:weak-maximum-nonpositive-zero-order} 和~\ref{cor:elliptic-comparison} 的对应结果也成立 (以 $D$ 代替 $\Omega$, 以 $\Sigma$ 代替 $\partial\Omega$). 注意, 令 $u=v\exp(\gamma t)$, 得到
\begin{equation}\label{eq:parabolic-zero-order-exponential-shift}
Lu=\exp(\gamma t)(Lv-\gamma v).
\end{equation}
因而若 $Lu\geq0$, 则 $(L-\gamma)v\geq0$, 反之亦然. 只要将 $\gamma$ 取得足够大, 总可以使 $c-\gamma\leq0$. 因此 Corollary~\ref{cor:elliptic-comparison} 的对应结果不需要对 $c$ 作任何符号假设.
